\documentclass[../main]{subfiles}
\ifSubfilesClassLoaded{\setcounter{chapter}{2}}{}
\begin{document}

\chapter{Pengenalan Unity Engine}

\begin{subcpmk}
  \item Sub-CPMK 3.1: Mengoperasikan panel Unity Editor dan mengelola Hierarchy serta Transform
\end{subcpmk}

% ============================================================
% MATERI POKOK
% ============================================================
\section{Sejarah dan Keunggulan Unity Engine}

\begin{learningoutcomes}
\item Mengetahui sejarah perkembangan Unity Engine
\item Memahami keunggulan Unity dibanding engine lain
\item Mengetahui platform yang didukung Unity
\item Memahami lisensi Unity
\end{learningoutcomes}

\subsection{Sejarah Unity}

Unity Technologies didirikan pada tahun 2004 di Denmark. Unity Engine pertama kali dirilis pada tahun 2005 dengan fokus pada pengembangan game untuk Mac. Sejak itu, Unity telah berkembang menjadi salah satu game engine paling populer di dunia \cite{UnityManual2024}.

\textbf{Timeline Penting}:
\begin{itemize}
    \item \textbf{2005}: Unity 1.0 dirilis (Mac only)
    \item \textbf{2009}: Unity 2.0 (Windows support)
    \item \textbf{2010}: Unity 3.0 (iOS \& Android support)
    \item \textbf{2012}: Unity 4.0 (WebGL, DirectX 11)
    \item \textbf{2017}: Unity 2017 (Timeline, Cinemachine)
    \item \textbf{2020}: Unity 2020 LTS (Stability focus)
    \item \textbf{2023}: Unity 6 (DOTS, Performance)
\end{itemize}

\subsection{Keunggulan Unity Engine}

\begin{enumerate}
\item \textbf{Cross-Platform}: Build sekali, deploy ke banyak platform
\item \textbf{User-Friendly}: Interface yang intuitif untuk pemula
\item \textbf{Asset Store}: Marketplace besar untuk assets dan tools
\item \textbf{Community}: Komunitas besar dan aktif
\item \textbf{Documentation}: Dokumentasi lengkap dan tutorial
\item \textbf{Free Version}: Unity Personal gratis untuk individu dan small teams
\item \textbf{C\# Scripting}: Bahasa pemrograman modern dan powerful
\item \textbf{Visual Editor}: WYSIWYG editor untuk level design
\end{enumerate}

\subsection{Platform yang Didukung}

Unity mendukung deployment ke berbagai platform: Desktop (Windows, Mac, Linux), Mobile (iOS, Android), Web (WebGL), Console (PlayStation, Xbox, Nintendo Switch), VR/AR (Oculus, HTC Vive, ARKit, ARCore).

\subsection{Lisensi Unity}

\textbf{Unity Personal}: Gratis untuk individu dan perusahaan dengan revenue < \$100K/tahun. \textbf{Unity Pro}: Berbayar dengan fitur tambahan. \textbf{Unity Enterprise}: Untuk perusahaan besar.

\begin{catatan}
Unity Personal sudah cukup untuk pembelajaran dan pengembangan game indie. Upgrade ke Pro hanya diperlukan jika membutuhkan fitur khusus atau untuk komersialisasi skala besar.
\end{catatan}

\section{Scene View dan Navigasi}

\subsection{Scene View}

Scene view memungkinkan Anda melihat dan mengedit dunia game secara visual. Tool navigasi (toolbar di atas Scene):

\begin{itemize}
  \item \textbf{Hand Tool (Q)}: Geser pandangan tanpa mengubah objek
  \item \textbf{Move (W)}: Pindahkan objek sepanjang sumbu X, Y, Z
  \item \textbf{Rotate (E)}: Putar objek
  \item \textbf{Scale (R)}: Ubah ukuran objek
  \item \textbf{Frame (F)}: Fokus kamera Scene pada objek terpilih
\end{itemize}

Navigasi orbit dan zoom dilakukan melalui kombinasi tombol Alt dengan klik mouse serta scroll.

\subsection{Konsep Project, Scene, dan Asset}

\textbf{Project} adalah kontainer untuk seluruh aset game (script, model, texture, material, scene). \textbf{Scene} adalah ruang kerja yang berisi GameObject dalam satu level atau layar. Project memiliki struktur folder (Scenes, Materials, Scripts, dll.) untuk mengorganisir aset. Scene disimpan sebagai file .unity di dalam project.

\subsection{Primitive 3D dan Material}

Unity menyediakan primitive 3D (Cube, Sphere, Plane, Cylinder, Capsule) sebagai geometri dasar. Setiap primitive memiliki Mesh Filter (data geometri) dan Mesh Renderer (menampilkan geometri ke layar). \textbf{Material} mendefinisikan tampilan permukaan (warna, tekstur); diterapkan ke Mesh Renderer. Material memisahkan tampilan dari geometri---satu material bisa dipakai banyak objek.

\subsection{Lighting}

\textbf{Directional Light} mensimulasikan sumber cahaya jauh (matahari). Properti Emission mengatur intensitas dan warna. Rotasi memengaruhi arah cahaya. Scene memerlukan minimal satu sumber cahaya agar objek 3D terlihat.

\section{Unity Editor Interface}

\begin{learningoutcomes}
\item Mengenal semua panel utama Unity Editor
\item Memahami fungsi setiap panel
\item Mampu mengatur layout editor
\item Mampu melakukan navigasi dasar di Unity Editor
\end{learningoutcomes}

\subsection{Panel Utama Unity Editor}

\begin{figure}[ht]
    \centering
    \includegraphics[width=\textwidth]{figures/unity_editor.png}
    \caption{Layout Standar Unity Editor}
    \label{fig:unity_editor}
\end{figure}

\subsubsection{Scene View}
Area untuk mengedit dan mengatur scene. Navigasi: Pan (Q), Rotate (E), Scale (R), Move (W). Gizmos untuk manipulasi objek.

\subsubsection{Game View}
Preview bagaimana game akan terlihat saat dimainkan. Aspect ratio dan resolution testing, Play mode preview.

\subsubsection{Hierarchy}
Menampilkan semua GameObject dalam scene aktif. Struktur parent-child, drag-and-drop untuk mengatur hierarchy.

\subsubsection{Inspector}
Menampilkan properti GameObject atau asset yang dipilih: Component properties, Transform values, Script variables.

\subsubsection{Project}
File browser untuk semua assets dalam project. Folder structure, asset preview, search dan filter.

\subsubsection{Console}
Menampilkan log, warning, dan error. Debug messages, error tracking.

\subsection{Shortcuts Penting}

\textbf{Ctrl/Cmd + S}: Save scene. \textbf{Ctrl/Cmd + Shift + N}: New GameObject. \textbf{F}: Focus pada objek terpilih. \textbf{Space}: Play/Pause. \textbf{Ctrl/Cmd + Shift + C}: Clear Console.

\section{Workflow Pengembangan Game di Unity}

\begin{learningoutcomes}
\item Memahami workflow dasar pengembangan game di Unity
\item Mengetahui urutan kerja yang efisien
\item Memahami konsep Scene dan Prefab
\item Mampu mengorganisir project dengan baik
\end{learningoutcomes}

\subsection{Workflow Dasar}

Planning, Asset Creation, Scene Setup, Scripting, Testing (Play mode), Iteration, Build untuk target platform.

\subsection{Konsep Scene}

\begin{definisi}[Scene]
Scene adalah container yang berisi semua GameObject, lighting, dan settings untuk level atau bagian tertentu dari game \cite{UnityManual2024}.
\end{definisi}

Best Practices: Satu scene per level, organize dengan folder dalam Hierarchy, gunakan empty GameObject sebagai organizer, save scene secara teratur.

\subsection{Konsep Prefab}

\begin{definisi}[Prefab]
Prefab adalah template GameObject yang dapat digunakan kembali dan diinstansiasi berkali-kali \cite{UnityManual2024}.
\end{definisi}

Gunakan Prefab untuk: GameObject yang digunakan berulang kali, konfigurasi kompleks, atau yang perlu diupdate secara konsisten.

\subsection{Organisasi Project}

Struktur folder yang disarankan: Scenes/, Scripts/, Art/, Audio/, Prefabs/, Materials/, Animations/.

\subsection{Version Control}

Unity bekerja dengan Git. Gunakan .gitignore untuk Unity. Jangan commit Library/ dan Temp/. Commit Assets/ dan ProjectSettings/.


% ============================================================
% AKTIVITAS PEMBELAJARAN
% ============================================================

\begin{aktivitas}
  \item \textbf{Setup Project}: Install Unity Hub dan Unity Editor (versi LTS), buat project baru dengan template 3D, eksplorasi struktur folder project.
  
  \item \textbf{Eksplorasi Interface}: Buka semua panel utama Unity Editor, coba setiap tool di toolbar, buat beberapa GameObject dan eksplorasi Hierarchy serta Inspector.
  
  \item \textbf{Setup Workflow}: Buat struktur folder yang terorganisir, buat scene baru dan setup basic hierarchy, buat beberapa GameObject dan convert ke Prefab.
\end{aktivitas}

% ============================================================
% LATIHAN DAN REFLEKSI
% ============================================================

\begin{latihan}
  \item Jelaskan sejarah singkat Unity Engine dan keunggulannya dibanding engine lain!
  
  \item Sebutkan langkah-langkah instalasi Unity Editor! Apa fungsi Unity Hub?
  
  \item Sebutkan 6 panel utama Unity Editor dan fungsinya masing-masing!
  
  \item Apa perbedaan antara Scene dan Prefab? Kapan sebaiknya menggunakan Prefab?
  
  \item Jelaskan workflow dasar pengembangan game di Unity!
  
  \item Sebutkan struktur folder project yang disarankan untuk project Unity!
  
  \item \textbf{Refleksi}: Apa tantangan yang Anda hadapi saat pertama kali menggunakan Unity Editor?
\end{latihan}

% ============================================================
% ASESMEN
% ============================================================

\begin{asesmen}
\textbf{Instrumen Penilaian untuk Sub-CPMK 3.1}

\textbf{A. Pilihan Ganda}

\begin{enumerate}
  \item Aplikasi untuk mengelola berbagai versi Unity dan project disebut:
  \begin{enumerate}
    \item Unity Editor
    \item Unity Hub
    \item Asset Store
    \item Package Manager
  \end{enumerate}
  
  \item Panel yang menampilkan semua GameObject dalam scene aktif adalah:
  \begin{enumerate}
    \item Inspector
    \item Project
    \item Hierarchy
    \item Console
  \end{enumerate}
  
  \item Template GameObject yang dapat digunakan kembali dan diinstansiasi berkali-kali disebut:
  \begin{enumerate}
    \item Scene
    \item Asset
    \item Prefab
    \item Component
  \end{enumerate}
\end{enumerate}

\textbf{B. Essay}

\begin{enumerate}
  \item Jelaskan workflow pengembangan game di Unity dari planning hingga build!
  \item Mengapa organisasi project penting? Sebutkan struktur folder yang disarankan!
\end{enumerate}

\textbf{Rubrik Penilaian}: Lihat Lampiran A
\end{asesmen}

% ============================================================
% CHECKLIST KOMPETENSI
% ============================================================

\begin{checklist}
  \item Saya dapat menjelaskan sejarah dan keunggulan Unity
  \item Saya dapat menginstal Unity Hub dan Unity Editor
  \item Saya dapat mengoperasikan panel Scene, Game, Hierarchy, Inspector, Project, Console
  \item Saya dapat membuat project baru dan memahami struktur foldernya
  \item Saya memahami konsep Scene dan Prefab
  \item Saya dapat mengelola Hierarchy dan Transform
\end{checklist}

% ============================================================
% RANGKUMAN
% ============================================================

\begin{rangkuman}
Bab ini membahas pengenalan Unity Engine termasuk sejarah, instalasi, interface editor, workflow, dan package management.

\textbf{Poin Kunci:}
\begin{itemize}
  \item Unity adalah game engine cross-platform yang powerful dan user-friendly
  \item Unity Hub untuk mengelola versi dan project
  \item Panel utama: Scene, Game, Hierarchy, Inspector, Project, Console
  \item Workflow: Planning, Asset Creation, Scene Setup, Scripting, Testing, Build
  \item Scene dan Prefab adalah konsep fundamental Unity
\end{itemize}
\end{rangkuman}

\ifSubfilesClassLoaded{
  \renewcommand{\bibname}{Daftar Pustaka}
  \bibliographystyle{plain}
  \bibliography{../references}
}{}
\end{document}
